\documentclass[12pt, a4paper]{article}
\usepackage[UTF8, zihao=-4]{ctex}
\usepackage[top=2.5cm, bottom=2.5cm, left=2.8cm, right=2.8cm]{geometry}
\usepackage{fontspec}
\setmainfont{Times New Roman}
\usepackage{setspace}
\setstretch{1.5}
\usepackage{booktabs}
\usepackage{array}
\usepackage{makecell}
\usepackage{multirow}
\usepackage{adjustbox}
\usepackage{amsmath, amssymb}
\usepackage{float}
\usepackage{caption}
\usepackage{fancyhdr}
\setlength{\headheight}{15pt}
\pagestyle{fancy}
\fancyhf{}
\fancyhead[C]{\small 实验数据记录单}
\fancyfoot[C]{\small\thepage}

\begin{document}

\begin{center}
  {\heiti\zihao{3} 马赫-曾德干涉实验 \\ 实验数据记录单} \\[1cm]
  \begin{tabular}{llll}
    姓\quad 名：& \underline{\hspace{4cm}} &
    学\quad 号：& \underline{\hspace{4cm}} \\[0.5cm]
    班\quad 级：& \underline{\hspace{4cm}} &
    实验日期：  & \underline{\hspace{4cm}} \\[0.5cm]
    指导教师：  & \underline{\hspace{4cm}} &
    教师签名：  & \underline{\hspace{4cm}} \\
  \end{tabular}
\end{center}
\vspace{1cm}

\subsection*{一、实验参数记录}
\begin{table}[H]
  \centering
  \begin{adjustbox}{max width=\textwidth}
    \begin{tabular}{cccc}
      \toprule
      激光波长 $\lambda$/nm & 气室长度 $L$ & 室温 $t$/$^\circ$C & 标准大气压 $p$ \\
      \midrule
      632.8 & & & \\
      \bottomrule
    \end{tabular}
  \end{adjustbox}
\end{table}
{\footnotesize 记录说明：除激光波长外，其余参数按现场仪器或教师要求填写；计算前需统一单位。}

\subsection*{二、光路调节检查}
\begin{table}[H]
  \centering
  \begin{adjustbox}{max width=\textwidth}
    \begin{tabular}{p{0.36\textwidth}p{0.52\textwidth}}
      \toprule
      检查项目 & 现场记录 \\
      \midrule
      激光近处、远处均过 2 mm 光阑中心 & \\
      空间滤波后光斑均匀，衍射环明显减弱或消失 & \\
      准直后近处、远处光斑大小基本不变 & \\
      两臂光程基本相等，光束高度一致 & \\
      合束后近处、远处两光斑均重合 & \\
      干涉条纹清晰且不过曝 & \\
      \bottomrule
    \end{tabular}
  \end{adjustbox}
\end{table}

\subsection*{三、光学平台稳定性观察}
\begin{table}[H]
  \centering
  \begin{adjustbox}{max width=\textwidth}
    \begin{tabular}{ccc}
      \toprule
      扰动方式 & 条纹变化现象 & 恢复情况 \\
      \midrule
      自然静置 & & \\
      轻拍平台 & & \\
      台旁地面跳动 & & \\
      \bottomrule
    \end{tabular}
  \end{adjustbox}
\end{table}

\subsection*{四、空气折射率测量数据}
\begin{table}[H]
  \centering
  \begin{adjustbox}{max width=\textwidth}
    \begin{tabular}{c*{9}{>{\centering\arraybackslash}p{1.35cm}}}
      \toprule
      序号 & 1 & 2 & 3 & 4 & 5 & 6 & 7 & 8 & 9 \\
      \midrule
      气压值/mmHg & & & & & & & & & \\
      条纹移动条数/个 & & & & & & & & & \\
      \bottomrule
    \end{tabular}
  \end{adjustbox}
\end{table}
{\footnotesize 记录说明：建议缓慢放气并记录累计条纹移动数；若记录相邻变化量，应在备注中说明。}

\subsection*{五、计算记录}
\noindent 采用公式：
\[
  n = 1 + \frac{\lambda}{L}\frac{\Delta K}{\Delta p}p
\]

\vspace{0.4cm}
\noindent 取用的数据点或拟合方法：

\vspace{1.2cm}

\noindent $\Delta K=$\underline{\hspace{3cm}}\quad
$\Delta p=$\underline{\hspace{3cm}}\quad
$\dfrac{\Delta K}{\Delta p}=$\underline{\hspace{3cm}}

\vspace{1cm}

\noindent 空气折射率实验值：$n=$\underline{\hspace{5cm}}

\vspace{0.8cm}

\noindent 与理论值比较及误差分析：

\vspace{2.2cm}

\noindent 备注：

\vspace{1.5cm}

\end{document}
